\documentclass[a4paper,12pt]{article}
\usepackage[utf8]{inputenc}
\usepackage[ngerman]{babel}
\usepackage{geometry}
\geometry{a4paper, left=3cm, right=2.5cm, top=2.5cm, bottom=2.5cm}

\usepackage[hidelinks]{hyperref}
\usepackage{enumitem}

\usepackage{longtable}
\usepackage{array}
\newcolumntype{L}[1]{>{\raggedright\arraybackslash}p{#1}}

\usepackage{graphicx}
\usepackage{float}
\usepackage[table]{xcolor}
\usepackage[most]{tcolorbox}


\usepackage{tabularx}




\title{Systementwurf\\[1ex]\large{Projektname: \textit{TippHub by THM}}}
\author{Gruppe C}
\date{\today}

\begin{document}
	
	\maketitle
	\newpage
	\tableofcontents
	\newpage
	
	
	\section{Operationelle Verträge}
	\subsection{LCB (Local Community Betreiber)}
		\newtcolorbox{contractbox}[1][]{
		colback=gray!5!white,
		colframe=gray!70!black,
		title=#1,
		fonttitle=\bfseries,
		boxrule=0.4pt,
		arc=4pt,
		outer arc=4pt,
		left=2mm,
		right=2mm,
		top=1mm,
		bottom=1mm,
		enhanced,
		breakable
	}
	
	\subsubsection{PUS in lokaler Community verwalten}
	\begin{contractbox}[\texttt{Operation OP1:openManagementTool}]
		\textbf{Operation:} openManagementTool()
		
			\textbf{Querverweis:} Use Cases "PUS in lokaler Community verwalten"
		
		\textbf{Vorbedingungen:}
		\begin{itemize}[noitemsep]
			\item LCB ist authentifiziert.
		\end{itemize}
		
		\textbf{Nachbedingungen:}
		\begin{itemize}[noitemsep]
			\item Eine Instanz von LcbAdministrationView lcb-administration wurde erzeugt.
			\item Das Verwaltungstool wurde geöffnet.
		\end{itemize}
	\end{contractbox}
	
	\begin{contractbox}[\texttt{Operation OP2:searchPUS}]
		\textbf{Operation:} searchPUS(username: String)
		
		\textbf{Querverweis:} Use Cases "PUS in lokaler Community verwalten"
		
		\textbf{Vorbedingungen:}
		\begin{itemize}[noitemsep]
			\item Das Verwaltungstool ist geöffnet.
		\end{itemize}
		
		\textbf{Nachbedingungen:}
		\begin{itemize}[noitemsep]
			\item lcb-administration wurde mit PeerUser assoziiert,
			basierend auf dem angegebenen username
		\end{itemize}
	\end{contractbox}
	
	\begin{contractbox}[\texttt{Operation OP3:displayPUS}]
	\textbf{Operation:} displayPUS()
	
	\textbf{Querverweis:} Use Cases "PUS in lokaler Community verwalten"
	
	\textbf{Vorbedingungen:}
	\begin{itemize}[noitemsep]
		\item Der LCB sucht nach einem PUS
	\end{itemize}
	
	\textbf{Nachbedingungen:}
	\begin{itemize}[noitemsep] 
		\item Der gesuchte PUS wurde erfolgreich gefunden.
	\end{itemize}
\end{contractbox}
	
	\begin{contractbox}[\texttt{Operation OP4:chooseAction}]
		\textbf{Operation:} chooseAction(action:Integer)
		
		\textbf{Querverweis:} Use Cases "PUS in lokaler Community verwalten"
		
		\textbf{Vorbedingungen:}
		\begin{itemize}[noitemsep]
			\item Der PUS wurde gefunden.
		\end{itemize}
		
		\textbf{Nachbedingungen:}
		\begin{itemize}[noitemsep]
			\item lcb-administration.action wird auf 1(Aufnahme) oder 2(Entfernung) gestellt.
		\end{itemize}
	\end{contractbox}
	
	\begin{contractbox}[\texttt{Operation OP5:confirmAction}]
		\textbf{Operation:} confirmAction(pusID:PusID, action:Integer)
		
		\textbf{Querverweis:} Use Cases "PUS in lokaler Community verwalten"
		
		\textbf{Vorbedingungen:}
		\begin{itemize}[noitemsep]
			\item Eine Aktion wurde gewählt.
		\end{itemize}
		
		\textbf{Nachbedingungen:}
		\begin{itemize}[noitemsep]
			\item Die lcb-administration.action wird dem System mitgeteilt
		\end{itemize}
	\end{contractbox}
	
	\begin{contractbox}[\texttt{Operation OP6:updateCommunityStatus}]
		\textbf{Operation:} updateCommunityStatus(pusID:PusID,localcommuityID:LocalcommuityID)
		
		\textbf{Querverweis:} Use Cases "PUS in lokaler Community verwalten"
		
		\textbf{Vorbedingungen:}
		\begin{itemize}[noitemsep]
			\item Die Aktion wurde bestätigt.
		\end{itemize}
		
		\textbf{Nachbedingungen:}
		\begin{itemize}[noitemsep]
			\item localcommunity wurde mit Peeruser assoziiert. 
			\item Der localcommunity.status des PUS wurde entsprechend der Aktion aktualisiert.
			\item Der PUS ist nun Mitglied oder kein Mitglied der Community.
		\end{itemize}
	\end{contractbox}
	
	
	\begin{contractbox}[\texttt{Operation OP7:notifySuccess}]
		\textbf{Operation:} notifySuccess()
		
		\textbf{Querverweis:} Use Cases "PUS in lokaler Community verwalten"
		
		\textbf{Vorbedingungen:}
		\begin{itemize}[noitemsep]
			\item Die Aktion wurde erfolgreich durchgeführt.
		\end{itemize}
		
		\textbf{Nachbedingungen:}
		\begin{itemize}[noitemsep]
			\item Eine Erfolgsmeldung über die durchgeführte Aktion wurde erzeugt.
		\end{itemize}
	\end{contractbox}
	
	\begin{contractbox}[\texttt{Operation OP8:displayFeedback}]
		\textbf{Operation:} displayFeedback(erfolg:boolean)
		
		\textbf{Querverweis:} Use Cases "PUS in lokaler Community verwalten"
		
		\textbf{Vorbedingungen:}
		\begin{itemize}[noitemsep]
			\item Die Aktion wurde erfolgreich durchgeführt.
		\end{itemize}
		
		\textbf{Nachbedingungen:}
		\begin{itemize}[noitemsep]
			\item Die Rückmeldung wird dem LCB angezeigt.
		\end{itemize}
	\end{contractbox}
	
	\subsubsection{Regionale Zugehörigkeit prüfen}
	\begin{contractbox}[\texttt{Operation OP1:openManagementTool}]
		\textbf{Operation:} openManagementTool()
		
		\textbf{Querverweis:} Use Cases "Regionale Zugehörigkeit prüfen"
		
		\textbf{Vorbedingungen:}
		\begin{itemize}[noitemsep]
			\item LCB ist authentifiziert.
		\end{itemize}
		
		\textbf{Nachbedingungen:}
		\begin{itemize}[noitemsep]
			\item Eine Instanz von LcbAdministrationView lcb-administration wurde erzeugt.
			\item Das Verwaltungstool wurde geöffnet.
		\end{itemize}
	\end{contractbox}
	
	\begin{contractbox}[\texttt{Operation OP2:searchPUS}]
		\textbf{Operation:} searchPUS(username: String)
		
		\textbf{Querverweis:} Use Cases "Regionale Zugehörigkeit prüfen"
		
		\textbf{Vorbedingungen:}
		\begin{itemize}[noitemsep]
			\item Das Verwaltungstool ist geöffnet.
		\end{itemize}
		
		\textbf{Nachbedingungen:}
		\begin{itemize}[noitemsep]
			\item lcb-administration wurde mit PeerUser assoziiert,
			basierend auf dem angegebenen username
		\end{itemize}
	\end{contractbox}
	
	\begin{contractbox}[\texttt{Operation OP3:displayPUS}]
		\textbf{Operation:} displayPUS()
		
		\textbf{Querverweis:} Use Cases "Regionale Zugehörigkeit prüfen"
		
		\textbf{Vorbedingungen:}
		\begin{itemize}[noitemsep]
			\item Der LCB sucht nach einem PUS
		\end{itemize}
		
		\textbf{Nachbedingungen:}
		\begin{itemize}[noitemsep] 
			\item Der gesuchte PUS wurde erfolgreich gefunden.
		\end{itemize}
	\end{contractbox}
	
	\begin{contractbox}[\texttt{Operation OP4:chooseAction}]
		\textbf{Operation:} chooseAction(action:Integer)
		
		\textbf{Querverweis:} Use Cases "Regionale Zugehörigkeit prüfen"
		
		\textbf{Vorbedingungen:}
		\begin{itemize}[noitemsep]
			\item Der PUS wurde gefunden.
		\end{itemize}
		
		\textbf{Nachbedingungen:}
		\begin{itemize}[noitemsep]
			\item lcb-administration.action wird auf 1(Aufnahme) oder 2(Entfernung) gestellt.
		\end{itemize}
	\end{contractbox}
	
	\begin{contractbox}[\texttt{Operation OP5:confirmAction}]
		\textbf{Operation:} confirmAction(pusID:PusID, action:Integer)
		
		\textbf{Querverweis:} Use Cases "Regionale Zugehörigkeit prüfen"
		
		\textbf{Vorbedingungen:}
		\begin{itemize}[noitemsep]
			\item Eine Aktion wurde gewählt.
		\end{itemize}
		
		\textbf{Nachbedingungen:}
		\begin{itemize}[noitemsep]
			\item Die lcb-administration.action wird dem System mitgeteilt
		\end{itemize}
	\end{contractbox}
	
	\begin{contractbox}[\texttt{Operation OP6:checkRegionalEligibility}]
		\textbf{Operation:} checkRegionalEligibility(pusID:PusID,localcommuityID:LocalcommuityID)
		
		\textbf{Querverweis:} Use Cases "Regionale Zugehörigkeit prüfen"
		
		\textbf{Vorbedingungen:}
		\begin{itemize}[noitemsep]
			\item Eine Aktion zur Aufnahme wurde bestätigt.
		\end{itemize}
		
		\textbf{Nachbedingungen:}
		\begin{itemize}[noitemsep]
			\item localcommunity wurde mit Peeruser assoziiert. 
			\item peeruser.region wird verglichen mit localcommunity.region.
			\item Das System hat überprüft, ob der PUS regional zur Community gehört.
		\end{itemize}
	\end{contractbox}
	
	
	\begin{contractbox}[\texttt{Operation OP7:updateCommunityStatus}]
		\textbf{Operation:} updateCommunityStatus(pusID:PusID,localcommuityID:LocalcommuityID)
		
		\textbf{Querverweis:} Use Cases "Regionale Zugehörigkeit prüfen"
		
		\textbf{Vorbedingungen:}
		\begin{itemize}[noitemsep]
			\item Die Aktion wurde bestätigt.
		\end{itemize}
		
		\textbf{Nachbedingungen:}
		\begin{itemize}[noitemsep]
			\item localcommunity wurde mit Peeruser assoziiert. 
			\item Der localcommunity.status des PUS wurde entsprechend der Aktion aktualisiert.
			\item Der PUS ist nun Mitglied oder kein Mitglied der Community.
		\end{itemize}
	\end{contractbox}
	
	
	\begin{contractbox}[\texttt{Operation OP8:notifySuccess}]
		\textbf{Operation:} notifySuccess()
		
		\textbf{Querverweis:} Use Cases "Regionale Zugehörigkeit prüfen"
		
		\textbf{Vorbedingungen:}
		\begin{itemize}[noitemsep]
			\item Die Aktion wurde erfolgreich durchgeführt.
		\end{itemize}
		
		\textbf{Nachbedingungen:}
		\begin{itemize}[noitemsep]
			\item Eine Erfolgsmeldung über die durchgeführte Aktion wurde erzeugt.
		\end{itemize}
	\end{contractbox}
	
	\begin{contractbox}[\texttt{Operation OP9:displayFeedback}]
		\textbf{Operation:} displayFeedback(erfolg:boolean)
		
		\textbf{Querverweis:} Use Cases "Regionale Zugehörigkeit prüfen"
		
		\textbf{Vorbedingungen:}
		\begin{itemize}[noitemsep]
			\item Die Aktion wurde erfolgreich durchgeführt.
		\end{itemize}
		
		\textbf{Nachbedingungen:}
		\begin{itemize}[noitemsep]
			\item Die Rückmeldung wird dem LCB angezeigt.
		\end{itemize}
	\end{contractbox}
	
	
	\subsubsection{Konfliktentscheidung zwischen PUSs auflösen}
	
	
	\begin{contractbox}[\texttt{Operation OP1:sendConflict}]
		\textbf{Operation:} sendConflict(pusID:List[PusID])
		
		\textbf{Querverweis:} Use Cases "Konfliktentscheidung zwischen PUSs auflösen"
		
		\textbf{Vorbedingungen:}
		\begin{itemize}[noitemsep]
			\item Ein Konflikt wurde erfasst.
		\end{itemize}
		
		\textbf{Nachbedingungen:}
		\begin{itemize}[noitemsep]
			\item Eine Instanz von Conflict conflict wurde erzeugt.
			\item conflict.status wurde auf 1 (offen) gesetzt.
			\item conflict.datum wurde auf das aktuelle Systemdatum gesetzt.
			\item confilct.details wurde mit dem übergebenen Text belegt.
			\item conflict wurde mit PeerUser assoziiert,
			basiert auf der pusID.
			\item Eine Benachrichtigung wurde an den Local Community Betreiber gesendet.
		\end{itemize}
	\end{contractbox}
	
	\begin{contractbox}[\texttt{Operation OP2:displayConflict}]
		\textbf{Operation:} displayConflict(conflictID:ConflictID)
		
		\textbf{Querverweis:} Use Cases "Konfliktentscheidung zwischen PUSs auflösen"
		
		\textbf{Vorbedingungen:}
		\begin{itemize}[noitemsep]
			\item Der Konflikt existiert.
		\end{itemize}
		
		\textbf{Nachbedingungen:}
		\begin{itemize}[noitemsep]
			\item Intern wurden Getter-Methoden auf dem Conflict Objekt aufgerufen
			\begin{itemize}
				\item \texttt{getStatus()}
				\item \texttt{getDate()}
				\item \texttt{getDetails()}
			\end{itemize}
			\item Die Daten des Konflikts werden angezeigt.
		\end{itemize}
	\end{contractbox}
	
	\begin{contractbox}[\texttt{Operation OP3:checkInvolvedPUS}]
		\textbf{Operation:} checkInvolvedPUS(conflictID:ConflictID)
		
		\textbf{Querverweis:} Use Cases "Konfliktentscheidung zwischen PUSs auflösen"
		
		\textbf{Vorbedingungen:}
		\begin{itemize}[noitemsep]
			\item Der Konflikt existiert.
		\end{itemize}
		
		\textbf{Nachbedingungen:}
		\begin{itemize}[noitemsep]
			\item Liste der beteiligten PUS wurde angezeigt.
		\end{itemize}
	\end{contractbox}
	
	\begin{contractbox}[\texttt{Operation OP4:contactPUS}]
		\textbf{Operation:} contactPUS(pusID:List[PusID])
		
		\textbf{Querverweis:} Use Cases "Konfliktentscheidung zwischen PUSs auflösen"
		
		\textbf{Vorbedingungen:}
		\begin{itemize}[noitemsep]
			\item Die PUSs sind dem Konflikt zugeordnet.
		\end{itemize}
		
		\textbf{Nachbedingungen:}
		\begin{itemize}[noitemsep]
			\item Die PUSs wurden kontaktiert.
		\end{itemize}
	\end{contractbox}
	

	\begin{contractbox}[\texttt{Operation OP5:submitDecision}]
		\textbf{Operation:} submitDecision(conflictID:ConflictID, decision:String)
		
	\textbf{Querverweis:} Use Cases "Konfliktentscheidung zwischen PUSs auflösen"
		
		\textbf{Vorbedingungen:}
		\begin{itemize}[noitemsep]
			\item Für den Konflikt existiert ein offener Status.
			\item Der LCB hat eine Entscheidung getroffen.
		\end{itemize}
		
		\textbf{Nachbedingungen:}
		\begin{itemize}[noitemsep]
			\item contract.decision wurde mit der übermittelten Entscheidung belegt.
		\end{itemize}
	\end{contractbox}
	
	\begin{contractbox}[\texttt{Operation OP6:updateConflictStatus}]
		\textbf{Operation:} updateConflictStatus(conflictID:ConflictID)
		
	\textbf{Querverweis:} Use Cases "Konfliktentscheidung zwischen PUSs auflösen"
		
		\textbf{Vorbedingungen:}
		\begin{itemize}[noitemsep]
			\item Eine gültige Entscheidung liegt vor.
		\end{itemize}
		
		\textbf{Nachbedingungen:}
		\begin{itemize}[noitemsep]
			\item conflict.status wurde auf 2(geschlossen) gesetzt.
		\end{itemize}
	\end{contractbox}
	
	
	
	\subsubsection{THMCoins an LCB gutschreiben}
	
	\begin{contractbox}[\texttt{Operation OP1:checkActiveLCBs()}]
		\textbf{Operation:} checkActiveLCBs()
		
			\textbf{Querverweis:} Use Cases "THMCoins an LCB gutschreiben"
		
		\textbf{Vorbedingungen:}
		\begin{itemize}[noitemsep]
			\item Der Systemprozess zur LCB-Aktivitätsprüfung ist aktiv.
		\end{itemize}
		
		\textbf{Nachbedingungen:}
		\begin{itemize}[noitemsep]
			\item localcommunity wurde mit localcommunitybetreiber assoziert.
			\item localcommuitybetreiber.status wurde überprüft
		\end{itemize}
	\end{contractbox}
	
	\begin{contractbox}[\texttt{Operation OP2:creditTHMCoins}]
		\textbf{Operation:} creditTHMCoins(pusID:PusID, amount:double)
		
		\textbf{Querverweis:} Use Cases "THMCoins an LCB gutschreiben"
		
		\textbf{Vorbedingungen:}
		\begin{itemize}[noitemsep]
			\item Der LCB ist aktiv.
		\end{itemize}
		
		\textbf{Nachbedingungen:}
		\begin{itemize}[noitemsep]
			\item peeruser wurde mit localcommunitybetreiber assoziert basierend auf der pusID.
			\item peeruser.thmcoins wurde um den angegebenen Betrag erhöht. 
			\item Eine Transaktion wurde durchgeführt, um dem LCB THMCoins gutzuschreiben.
		\end{itemize}
	\end{contractbox}
	
	
	\begin{contractbox}[\texttt{Operation OP3:notifyLCB}]
		\textbf{Operation:} notifyLCB(pusID:PusID)
		
		\textbf{Querverweis:} Use Case "THMCoins an LCB gutschreiben"
		
		\textbf{Vorbedingungen:}
		\begin{itemize}[noitemsep]
			\item Eine erfolgreiche Gutschrift-Transaktion wurde durchgeführt.
		\end{itemize}
		
		\textbf{Nachbedingungen:}
		\begin{itemize}[noitemsep]
			\item Der LCB wurde über die Gutschrift informiert.
			\item Eine Benachrichtigung mit dem übergebenen Inhalt wurde dem Benutzer angezeigt.
		\end{itemize}
	\end{contractbox}
	
	\subsubsection{Lokales Community-Event anlegen}
	
	\begin{contractbox}[\texttt{Operation OP1:openEventTool}]
		\textbf{Operation:} openEventTool()
		
		\textbf{Querverweis:} Use Cases "Lokales Community-Event anlegen"
		
		\textbf{Vorbedingungen:}
		\begin{itemize}[noitemsep]
			\item LCB ist authentifiziert.
		\end{itemize}
		
		\textbf{Nachbedingungen:}
		\begin{itemize}[noitemsep]
			\item Eine Instanz von EventManagementView  eventmanagementview wurde erzeugt.
			\item Das Event-Verwaltungstool wurde geöffnet.
		\end{itemize}
	\end{contractbox}
	
	\begin{contractbox}[\texttt{Operation OP2:enterEventDetails}]
		\textbf{Operation:} enterEventDetails()
		
		\textbf{Querverweis:} Use Cases "Lokales Community-Event anlegen"
		
		\textbf{Vorbedingungen:}
		\begin{itemize}[noitemsep]
			\item Das Event-Verwaltungstool ist geöffnet.
		\end{itemize}
		
		\textbf{Nachbedingungen:}
		\begin{itemize}[noitemsep] 
			\item Die eingegebenen Eventdetails wurden eingetippt.
		\end{itemize}
	\end{contractbox}
	
	
	
	\begin{contractbox}[\texttt{Operation OP3:saveEvent}]
		\textbf{Operation:} saveEvent()
		
		\textbf{Querverweis:} Use Cases "Lokales Community-Event anlegen"
		
		\textbf{Vorbedingungen:}
		\begin{itemize}[noitemsep]
			\item Gültige Eventdetails wurden eingegeben.
		\end{itemize}
		
		\textbf{Nachbedingungen:}
		\begin{itemize}[noitemsep]
			\item Intern wurden Setter-Methoden auf dem eventmanagementview Objekt aufgerufen
			\begin{itemize}[noitemsep]
				\item \texttt{setTitle()}
				\item \texttt{setDate()}
				\item \texttt{setDescription()}
			\end{itemize}
			\item Die eingegebenen Eventdetails wurden validiert.
		\end{itemize}
	\end{contractbox}
	
	\begin{contractbox}[\texttt{Operation OP4:displayFeedback}]
		\textbf{Operation:} displayFeedback(erfolg:boolean)
		
		\textbf{Querverweis:} Use Cases "Lokales Community-Event anlegen"
		
		\textbf{Vorbedingungen:}
		\begin{itemize}[noitemsep]
			\item Die Aktion wurde erfolgreich durchgeführt.
		\end{itemize}
		
		\textbf{Nachbedingungen:}
		\begin{itemize}[noitemsep]
			\item Die Rückmeldung wird dem LCB angezeigt.
		\end{itemize}
	\end{contractbox}
	
	
	\begin{contractbox}[\texttt{Operation OP5:assignEventToLocalCommunity}]
		\textbf{Operation:} assignEventToLocalCommunity(eventID:EventID, localcommunityID:LocalCommunityID)
		
		\textbf{Querverweis:} Use Cases "Lokales Community-Event anlegen"
		
		\textbf{Vorbedingungen:}
		\begin{itemize}[noitemsep]
			\item Das Event wurde erfolgreich angelegt.
			\item Die Ziel-Community existiert.
		\end{itemize}
		
		\textbf{Nachbedingungen:}
		\begin{itemize}[noitemsep]
			\item eventmanagementview wurde mit localcommunity assoziiert.
			\item Das Event wurde der lokalen Community des LCB zugeordnet.
		\end{itemize}
	\end{contractbox}
	
	\begin{contractbox}[\texttt{Operation OP6:notifyPUSAboutEvent}]
		\textbf{Operation:} notifyPUSAboutEvent(eventID:EventID)
		
		\textbf{Querverweis:} Use Cases "Lokales Community-Event anlegen"
		
		\textbf{Vorbedingungen:}
		\begin{itemize}[noitemsep]
			\item Das Event wurde einer lokalen Community zugeordnet.
			\item Es existieren PUSs in dieser Community.
		\end{itemize}
		
		\textbf{Nachbedingungen:}
		\begin{itemize}[noitemsep]
			\item eventmanagementview wurde mit localcommunity assoziiert.
			\item localcommunity wurde mit peeruser assoziiert.
			\item Eine Benachrichtigung über das Event wurde an alle relevanten PUSs gesendet.
		\end{itemize}
	\end{contractbox}

	\subsection{LCB (TippHub Betreiber)}
		\newtcolorbox{contractbox1}[1][]{
		colback=gray!5!white,
		colframe=gray!70!black,
		title=#1,
		fonttitle=\bfseries,
		boxrule=0.4pt,
		arc=4pt,
		outer arc=4pt,
		left=2mm,
		right=2mm,
		top=1mm,
		bottom=1mm,
		enhanced,
		breakable
	}

	\subsubsection{Coinkaufpreis ändern}
	\begin{contractbox}[\texttt{Operation OP1(UC-1): openCoinManagementTool}]
		\textbf{Operation:} openCoinManagementTool()
		
			\textbf{Querverweis:} UC-1: Coinkaufpreis ändern
		
		\textbf{Vorbedingungen:}
		\begin{itemize}[noitemsep]
			\item THB ist authentifiziert.
		\end{itemize}
		
		\textbf{Nachbedingungen:}
		\begin{itemize}[noitemsep]
			\item Eine Instanz von ThbAdministrationView thb-administration wurde erzeugt.
			\item Das Coin-Verwaltungstool wurde geöffnet.
		\end{itemize}
	\end{contractbox}

	\begin{contractbox}[\texttt{Operation OP2(UC-1): viewCoinPrice}]
		\textbf{Operation:} viewCoinPrice()
		
			\textbf{Querverweis:} UC-1: Coinkaufpreis ändern
		
		\textbf{Vorbedingungen:}
		\begin{itemize}[noitemsep]
			\item THB ist authentifiziert.
		\end{itemize}
		
		\textbf{Nachbedingungen:}
		\begin{itemize}[noitemsep]
			\item Intern wurde Getter-Methoden auf dem Coin Objekt aufgerufen: getCoinPrice()
			\item Die View des Coin-Verwaltungstools wurde angezeigt.
		\end{itemize}
	\end{contractbox}
	
	\begin{contractbox}[\texttt{Operation OP3(UC-1): setCoinPrice}]
		\textbf{Operation:} setCoinPrice(price: Float)
		
		\textbf{Querverweis:} UC-1: Coinkaufpreis ändern
		
		\textbf{Vorbedingungen:}
		\begin{itemize}[noitemsep]
			\item Das Coin-Verwaltungstool ist geöffnet.
			\item Mit getCoinPrice() wird der aktuelle Preis einer ThmCoin zurückgegeben.
		\end{itemize}
		
		\textbf{Nachbedingungen:}
		\begin{itemize}[noitemsep]
			\item Der Werte von thmcoin.price, wurde auf den Eingegebenen Wert gesetzt.
		\end{itemize}
	\end{contractbox}

	\begin{contractbox}[\texttt{Operation OP4(UC-1): showCoinChangeConfirmation}]
		\textbf{Operation:} showCoinChangeConfirmation(coinprice)
		
			\textbf{Querverweis:} UC-1: Coinkaufpreis ändern
		
		\textbf{Vorbedingungen:}
		\begin{itemize}[noitemsep]
			\item Der Coinpreis wurde geändert
		\end{itemize}
		
		\textbf{Nachbedingungen:}
		\begin{itemize}[noitemsep]
			\item Es wurde dem Nutzer mit einer Popup-Anzeige, das Anlegen der neuen Community bestätigt.
		\end{itemize}
	\end{contractbox}

	\subsubsection{User sperren}
	\begin{contractbox}[\texttt{Operation OP1(UC-2): openUserManagementTool}]
		\textbf{Operation:} openUserManagementTool()
		
			\textbf{Querverweis:} UC-2: User sperren
		
		\textbf{Vorbedingungen:}
		\begin{itemize}[noitemsep]
			\item THB ist authentifiziert.
		\end{itemize}
		
		\textbf{Nachbedingungen:}
		\begin{itemize}[noitemsep]
			\item Eine Instanz von ThbAdministrationView thb-administration wurde erzeugt.
			\item Das User-Verwaltungstool wurde geöffnet.
		\end{itemize}
	\end{contractbox}

	\begin{contractbox}[\texttt{Operation OP2(UC-2): viewUserList}]
		\textbf{Operation:} viewUserList()
		
			\textbf{Querverweis:} UC-2: User sperren
		
		\textbf{Vorbedingungen:}
		\begin{itemize}[noitemsep]
			\item THB ist authentifiziert.
		\end{itemize}
		
		\textbf{Nachbedingungen:}
		\begin{itemize}[noitemsep]
			\item Mit getUserDetails(), wird eine Liste der User und ihr jeweiliger Status zurückgegeben.
			\item Die View des User-Verwaltungstools wurde angezeigt.
		\end{itemize}
	\end{contractbox}
	
	\begin{contractbox}[\texttt{Operation OP3(UC-2): setUserStatus}]
		\textbf{Operation:} setUserStatus(status: Int)
		
		\textbf{Querverweis:} UC-2: User sperren
		
		\textbf{Vorbedingungen:}
		\begin{itemize}[noitemsep]
			\item Das User-Verwaltungstool ist geöffnet.
		\end{itemize}
		
		\textbf{Nachbedingungen:}
		\begin{itemize}[noitemsep]
			\item Das Attribut, User.status, wurde auf 1 gesetzt.
		\end{itemize}
	\end{contractbox}

	\begin{contractbox}[\texttt{Operation OP4(UC-2): showStatusChangeConfirmation}]
		\textbf{Operation:} showStatusChangeConfirmation(userid: Int, status: Int)
		
			\textbf{Querverweis:} UC-2: User sperren
		
		\textbf{Vorbedingungen:}
		\begin{itemize}[noitemsep]
			\item Der Userstatus wurde geändert.
		\end{itemize}
		
		\textbf{Nachbedingungen:}
		\begin{itemize}[noitemsep]
			\item Es wurde dem Nutzer mit einer Popup-Anzeige, die Statusänderung eines Users bestätigt.
		\end{itemize}
	\end{contractbox}
	
	\subsubsection{User freigeben}
	\begin{contractbox}[\texttt{Operation OP1(UC-3): openUserManagementTool}]
		\textbf{Operation:} openUserManagementTool()
		
			\textbf{Querverweis:} UC-3: User freigeben
		
		\textbf{Vorbedingungen:}
		\begin{itemize}[noitemsep]
			\item THB ist authentifiziert.
		\end{itemize}
		
		\textbf{Nachbedingungen:}
		\begin{itemize}[noitemsep]
			\item Eine Instanz von ThbAdministrationView thb-administration wurde erzeugt.
			\item Das User-Verwaltungstool wurde geöffnet.
			\item Mit getUserDetails(), wird eine Liste der User und ihr jeweiliger Status zurückgegeben.
		\end{itemize}
	\end{contractbox}

	\begin{contractbox}[\texttt{Operation OP2(UC-3): viewUserList}]
		\textbf{Operation:} viewUserList()
		
			\textbf{Querverweis:} UC-3: User freigeben
		
		\textbf{Vorbedingungen:}
		\begin{itemize}[noitemsep]
			\item THB ist authentifiziert.
		\end{itemize}
		
		\textbf{Nachbedingungen:}
		\begin{itemize}[noitemsep]
			\item Mit getUserDetails(), wird eine Liste der User und ihr jeweiliger Status zurückgegeben.
			\item Die View des User-Verwaltungstools wurde angezeigt.
		\end{itemize}
	\end{contractbox}
	
	\begin{contractbox}[\texttt{Operation OP3(UC-3): setUserStatus}]
		\textbf{Operation:} setUserStatus(status: Float)
		
		\textbf{Querverweis:} UC-3: User freigeben
		
		\textbf{Vorbedingungen:}
		\begin{itemize}[noitemsep]
			\item Das User-Verwaltungstool ist geöffnet.
		\end{itemize}
		
		\textbf{Nachbedingungen:}
		\begin{itemize}[noitemsep]
			\item Das Attribut, User.status, wurde auf 0 gesetzt.
		\end{itemize}
	\end{contractbox}

	\begin{contractbox}[\texttt{Operation OP4(UC-3): showStatusChangeConfirmation}]
		\textbf{Operation:} showStatusChangeConfirmation(userid: Int, status: Int)
		
			\textbf{Querverweis:} UC-3: User freigeben
		
		\textbf{Vorbedingungen:}
		\begin{itemize}[noitemsep]
			\item Der Userstatus wurde geändert.
		\end{itemize}
		
		\textbf{Nachbedingungen:}
		\begin{itemize}[noitemsep]
			\item Es wurde dem Nutzer mit einer Popup-Anzeige, die Statusänderung eines Users bestätigt.
		\end{itemize}
	\end{contractbox}

	\subsubsection{Kennzahlen der Plattform ansehen}
	\begin{contractbox}[\texttt{Operation OP1(UC-4): openDashboard}]
		\textbf{Operation:} openDashboard()
		
			\textbf{Querverweis:} UC-4: Kennzahlen der Plattform ansehen
		
		\textbf{Vorbedingungen:}
		\begin{itemize}[noitemsep]
			\item THB ist authentifiziert.
		\end{itemize}
		
		\textbf{Nachbedingungen:}
		\begin{itemize}[noitemsep]
			\item Eine Instanz von ThbAdministrationView thb-administration wurde erzeugt.
			\item Das Dashboard-Verwaltungstool wurde geöffnet.
			\item Mit getDashboardPreview(), wird ein Überblick der Kennzahlen der letzten Woche zurückgegeben.
		\end{itemize}
	\end{contractbox}

	\begin{contractbox}[\texttt{Operation OP2(UC-4): viewDashboardPreview}]
		\textbf{Operation:} viewDashboardPreview()
		
			\textbf{Querverweis:} UC-4: Kennzahlen der Plattform ansehen
		
		\textbf{Vorbedingungen:}
		\begin{itemize}[noitemsep]
			\item Die Dashboardview wurde erzeugt.
		\end{itemize}
		
		\textbf{Nachbedingungen:}
		\begin{itemize}[noitemsep]
			\item Mit getDashboardPreview(), wird eine Liste der User und ihr jeweiliger Status zurückgegeben.
			\item Die View des User-Verwaltungstools wurde angezeigt.
		\end{itemize}
	\end{contractbox}

	\subsubsection{Details der Kennzahlen ansehen}
	\begin{contractbox}[\texttt{Operation OP1(UC-5): openDashboard}]
		\textbf{Operation:} openDashboard()
		
			\textbf{Querverweis:} UC-5: Kennzahlen der Plattform ansehen
		
		\textbf{Vorbedingungen:}
		\begin{itemize}[noitemsep]
			\item THB ist authentifiziert.
		\end{itemize}
		
		\textbf{Nachbedingungen:}
		\begin{itemize}[noitemsep]
			\item Eine Instanz von ThbAdministrationView thb-administration wurde erzeugt.
			\item Das Dashboard-Verwaltungstool wurde geöffnet.
			\item Mit getDashboardPreview(), wird ein Überblick der Kennzahlen der letzten Woche zurückgegeben.
		\end{itemize}
	\end{contractbox}

	\begin{contractbox}[\texttt{Operation OP2(UC-5): viewDashboardPreview}]
		\textbf{Operation:} viewDashboardPreview()
		
			\textbf{Querverweis:} UC-5: Kennzahlen der Plattform ansehen
		
		\textbf{Vorbedingungen:}
		\begin{itemize}[noitemsep]
			\item Die Dashboardview wurde erzeugt.
		\end{itemize}
		
		\textbf{Nachbedingungen:}
		\begin{itemize}[noitemsep]
			\item Mit getDashboardPreview(), wird eine Liste der User und ihr jeweiliger Status zurückgegeben.
			\item Die View des User-Verwaltungstools wurde angezeigt.
		\end{itemize}
	\end{contractbox}

	\begin{contractbox}[\texttt{Operation OP3(UC-5): showDetailedNumbers}]
		\textbf{Operation:} showDetailedNumbers()
		
			\textbf{Querverweis:} UC-5: Kennzahlen der Plattform ansehen
		
		\textbf{Vorbedingungen:}
		\begin{itemize}[noitemsep]
			\item OP1(openDashboard()) wurde ausgeführt.
		\end{itemize}
		
		\textbf{Nachbedingungen:}
		\begin{itemize}[noitemsep]
			\item Mit getRegPusDetails(), getLocalComDetails(), getBettingActDetails(), oder getCoinDetails() werden jeweils Detailierte Kennzahlen zurückgeben.
			\item Dashboard wurde mit Peer User assoziiert.
			\item Dashborad wurde mit Local Community assoziiert.
			\item Dashborad wurde mit Wettanfrage assoziiert.
			\item Dashborad wurde mit THM Coin assoziiert.
		\end{itemize}
	\end{contractbox}

	\begin{contractbox}[\texttt{Operation OP4(UC-5): showDetailedNumbers(Daterange)}]
		\textbf{Operation:} showDetailedNumbers(Daterange)
		
			\textbf{Querverweis:} UC-5: Kennzahlen der Plattform ansehen
		
		\textbf{Vorbedingungen:}
		\begin{itemize}[noitemsep]
			\item OP1(openDashboard()) wurde ausgeführt.
		\end{itemize}
		
		\textbf{Nachbedingungen:}
		\begin{itemize}[noitemsep]
			\item Mit getRegPusDetails(Daterange), getLocalComDetails(Daterange), getBettingActDetails(Daterange), oder getCoinDetails(Daterange) wurden Detailierte Kennzahlen zurückgeben.
			\item Dashboard wurde mit Peer User assoziiert.
			\item Dashborad wurde mit Local Community assoziiert.
			\item Dashborad wurde mit Wettanfrage assoziiert.
			\item Dashborad wurde mit THM Coin assoziiert
		\end{itemize}
	\end{contractbox}

	\subsubsection{Lokale Community anlegen}
	\begin{contractbox}[\texttt{Operation OP1(UC-6): OpenCommunityTool}]
		\textbf{Operation:} openCommunityTool()
		
			\textbf{Querverweis:} UC-6: Lokale Community anlegen
		
		\textbf{Vorbedingungen:}
		\begin{itemize}[noitemsep]
			\item THB ist authentifiziert.
		\end{itemize}
		
		\textbf{Nachbedingungen:}
		\begin{itemize}[noitemsep]
			\item Eine Instanz von ThbCommunityView thb-community wurde erzeugt.
			\item Das Community-Verwaltungstool wurde geöffnet.
			\item Mit getCommunityList() wird eine Auflistung der bereits existierenden lokalen Communities angezeigt.
		\end{itemize}
	\end{contractbox}

	\begin{contractbox}[\texttt{Operation OP2(UC-6): viewCommunityform}]
		\textbf{Operation:} viewCommunityForm()
		
			\textbf{Querverweis:} UC-6: Lokale Community anlegen
		
		\textbf{Vorbedingungen:}
		\begin{itemize}[noitemsep]
			\item OP1(openCommunityTool()) wurde ausgeführt.
		\end{itemize}
		
		\textbf{Nachbedingungen:}
		\begin{itemize}[noitemsep]
			\item Mit viewCommunityForm() wurde der Server zur Anzeige eines neuen Community-Formulars aufgefordert.
			\item Mit getCommunities() wurden die notwendigen Daten für das Frontend angefordert.
			\item Das Formular wird dem Nutzer angezeigt.
		\end{itemize}
	\end{contractbox}

	\begin{contractbox}[\texttt{Operation OP3(UC-6): Send Community-Form}]
		\textbf{Operation:} sendCommunityForm(communityName: String, lcbName)
		
			\textbf{Querverweis:} UC-6: Lokale Community anlegen
		
		\textbf{Vorbedingungen:}
		\begin{itemize}[noitemsep]
			\item OP1(openCommunityTool()) wurde ausgeführt.
			\item OP2(viewCommunityForm()) wurde ausgeführt.
		\end{itemize}
		
		\textbf{Nachbedingungen:}
		\begin{itemize}[noitemsep]
			\item Mit sendCommunityForm(communityName, managerName) wurden die Daten an den Server gesendet.
			\item Es wurde mit newCommunity() ein neues Objekt von localCommunity angelegt.
			\item Es wurde mit setCommunityName(communityName) der Name der Community gesetzt.
			\item Es wurde mit setCommunityManager(lcbName) der zugehörige Local Community Betreiber gesetzt.
		\end{itemize}
	\end{contractbox}

	\begin{contractbox}[\texttt{Operation OP4(UC-6): Confirm new Community}]
		\textbf{Operation:} showNewCommunityConfirmation(communityName, lcbName)
		
			\textbf{Querverweis:} UC-6: Lokale Community anlegen
		
		\textbf{Vorbedingungen:}
		\begin{itemize}[noitemsep]
			\item OP1(openCommunityTool()) wurde ausgeführt.
			\item OP2(viewCommunityForm()) wurde ausgeführt.
			\item OP3(sendCommunityForm(communityName, lcbName)) wurde ausgeführt.
		\end{itemize}
		
		\textbf{Nachbedingungen:}
		\begin{itemize}[noitemsep]
			\item Es wurde dem Nutzer mit einer Popup-Anzeige, das Anlegen der neuen Community bestätigt.
		\end{itemize}
	\end{contractbox}

	\subsubsection{Lokale Community löschen}
	\begin{contractbox}[\texttt{Operation OP1(UC-7): OpenCommunityTool}]
		\textbf{Operation:} openCommunityTool()
		
			\textbf{Querverweis:} UC-7: Lokale Community löschen
		
		\textbf{Vorbedingungen:}
		\begin{itemize}[noitemsep]
			\item THB ist authentifiziert.
		\end{itemize}
		
		\textbf{Nachbedingungen:}
		\begin{itemize}[noitemsep]
			\item Eine Instanz von ThbCommunityView thb-community wurde erzeugt.
			\item Das Community-Verwaltungstool wurde geöffnet.
		\end{itemize}
	\end{contractbox}

	\begin{contractbox}[\texttt{Operation OP2(UC-7): viewCommunityList}]
		\textbf{Operation:} viewCommunityList()
		
			\textbf{Querverweis:} UC-7: Lokale Community anlegen
		
		\textbf{Vorbedingungen:}
		\begin{itemize}[noitemsep]
			\item OP1(openCommunityTool()) wurde ausgeführt.
		\end{itemize}
		
		\textbf{Nachbedingungen:}
		\begin{itemize}[noitemsep]
			\item Mit getCommunities() wurden die notwendigen Daten für das Frontend angefordert und eine Liste der bestehenden Communities anzuzeigen.
		\end{itemize}
	\end{contractbox}

	\begin{contractbox}[\texttt{Operation OP3(UC-7): delete Community}]
		\textbf{Operation:} deleteCommunity(communityName)
		
			\textbf{Querverweis:} UC-7: Lokale Community löschen
		
		\textbf{Vorbedingungen:}
		\begin{itemize}[noitemsep]
			\item OP1(openCommunityTool()) wurde ausgeführt.
		\end{itemize}
		
		\textbf{Nachbedingungen:}
		\begin{itemize}[noitemsep]
			\item Mit deleteCommunity(communityName) wurde eine Aufforderung an den Server geschickt, ein Community zu löschen.
			\item Der Server prüft die Eingabe auf Korrektheit.
		\end{itemize}
	\end{contractbox}

	\begin{contractbox}[\texttt{Operation OP4(UC-7): Confirm Community deletion}]
		\textbf{Operation:} confirmCommunityDeletion(communityName)
		
			\textbf{Querverweis:} UC-7: Lokale Community löschen
		
		\textbf{Vorbedingungen:}
		\begin{itemize}[noitemsep]
			\item OP1(openCommunityTool()) wurde ausgeführt.
			\item OP2(deleteCommunity(communityName)) wurde ausgeführt.
		\end{itemize}
		
		\textbf{Nachbedingungen:}
		\begin{itemize}[noitemsep]
			\item Der Nutzer wurde gefragt, ob er die Community wirklich löschen möchte.
		\end{itemize}
	\end{contractbox}

	\begin{contractbox}[\texttt{Operation OP5(UC-7): Community wird gelöscht}]
		\textbf{Operation:} deleteCommunity(communityName)
		
			\textbf{Querverweis:} UC-7: Lokale Community löschen
		
		\textbf{Vorbedingungen:}
		\begin{itemize}[noitemsep]
			\item OP1(openCommunityTool()) wurde ausgeführt.
			\item OP2(viewCommunityForm()) wurde ausgeführt.
			\item OP3(confirmCommunityDeletion(communityName)) wurde bestätigt.
		\end{itemize}
		
		\textbf{Nachbedingungen:}
		\begin{itemize}[noitemsep]
			\item Mit deleteCommunity(communityName), wurde das Objekt zugehörige Objekt von localCommunity gelöscht.
		\end{itemize}
	\end{contractbox}

	\begin{contractbox}[\texttt{Operation OP6(UC-7): Confirm Community Deletion}]
		\textbf{Operation:} showNewCommunityDeletionConfirmation(communityName)
		
			\textbf{Querverweis:} UC-7: Lokale Community löschen
		
		\textbf{Vorbedingungen:}
		\begin{itemize}[noitemsep]
			\item OP1(openCommunityTool()) wurde ausgeführt.
			\item OP2(viewCommunityForm()) wurde ausgeführt.
			\item OP3(confirmCommunityDeletion(communityName)) wurde bestätigt.
			\item Op4(deleteCommunity(communityName)) wurde ausgeführt.
		\end{itemize}
		
		\textbf{Nachbedingungen:}
		\begin{itemize}[noitemsep]
			\item Es wurde dem Nutzer mit einer Popup-Anzeige, das Löschen der Community bestätigt.
		\end{itemize}
	\end{contractbox}

	\subsubsection{Community-Manager zuordnen }
	\begin{contractbox}[\texttt{Operation OP1(UC-8): OpenCommunityTool}]
		\textbf{Operation:} openCommunityTool()
		
			\textbf{Querverweis:} UC-8: Community-Manager zuordnen
		
		\textbf{Vorbedingungen:}
		\begin{itemize}[noitemsep]
			\item THB ist authentifiziert.
		\end{itemize}
		
		\textbf{Nachbedingungen:}
		\begin{itemize}[noitemsep]
			\item Eine Instanz von ThbCommunityView thb-community wurde erzeugt.
			\item Das Community-Verwaltungstool wurde geöffnet.
		\end{itemize}
	\end{contractbox}

	\begin{contractbox}[\texttt{Operation OP2(UC-8): viewCommunityList}]
		\textbf{Operation:} viewCommunityList()
		
			\textbf{Querverweis:} UC-8: Community-Manager zuordnen
		
		\textbf{Vorbedingungen:}
		\begin{itemize}[noitemsep]
			\item OP1(openCommunityTool()) wurde ausgeführt.
		\end{itemize}
		
		\textbf{Nachbedingungen:}
		\begin{itemize}[noitemsep]
			\item Mit getCommunities() wurden die notwendigen Daten für das Frontend angefordert und eine Liste der bestehenden Communities anzuzeigen.
		\end{itemize}
	\end{contractbox}

	\begin{contractbox}[\texttt{Operation OP3(UC-8): Edit Community}]
		\textbf{Operation:} editCommunity(communityName)
		
			\textbf{Querverweis:} UC-8: Community-Manager zuordnen
		
		\textbf{Vorbedingungen:}
		\begin{itemize}[noitemsep]
			\item OP1(openCommunityTool()) wurde ausgeführt.
		\end{itemize}
		
		\textbf{Nachbedingungen:}
		\begin{itemize}[noitemsep]
			\item Mit editCommunity(communityName) wurde eine Detailansicht der lokalen Community geöffnet.
		\end{itemize}
	\end{contractbox}

	\begin{contractbox}[\texttt{Operation OP4(UC-8): Change Community-Manager}]
		\textbf{Operation:} changeCommunityManger(communityName, PusID)
		
			\textbf{Querverweis:} UC-8: Community-Manager zuordnen
			
		\textbf{Vorbedingungen:}
		\begin{itemize}[noitemsep]
			\item OP1(openCommunityTool()) wurde ausgeführt.
			\item OP2(editCommunity(communityName)) wurde ausgeführt.
		\end{itemize}
		
		\textbf{Nachbedingungen:}
		\begin{itemize}[noitemsep]
			\item In Local Community, wurde die Variable Community-Manager, auf die in der Methode übergebeben PusID gesetzt.
		\end{itemize}
	\end{contractbox}

	\begin{contractbox}[\texttt{Operation OP5(UC-8): Confirm Manager-Change}]
		\textbf{Operation:} showNewMangerChangeConfirmation(communityName, PusID)
		
			\textbf{Querverweis:} C-8: Community-Manager zuordnen
		
		\textbf{Vorbedingungen:}
		\begin{itemize}[noitemsep]
			\item OP1(openCommunityTool()) wurde ausgeführt.
			\item OP2(editCommunity(communityName)) wurde ausgeführt.
			\item OP3(changeCommunityManger(communityName, PusID)) wurde erfolgreich ausgeführt.
		\end{itemize}
		
		\textbf{Nachbedingungen:}
		\begin{itemize}[noitemsep]
			\item Es wurde dem Nutzer mit einer Popup-Anzeige, die Änderunbvg des Community-Managers bestätigt.
		\end{itemize}
	\end{contractbox}




\end{document}